% chapters/tikz_ch09.tex - 第 9 章全部插图的高精度原生 TikZ 矢量代码

% 1. 图 9.1: 切空间 T_pM 中的光锥
\newcommand{\drawFigNineOne}{%
\begin{tikzpicture}[scale=1.05, >=Stealth, line join=round]
    % 坐标轴
    \draw[->] (0, -2.55) -- (0, 3.05) node[above] {$t$};
    \draw[->] (-2.85, 0) -- (3.05, 0) node[right] {$x$};
    \draw[->] (0, 0) -- (-1.75, -1.35) node[below left] {$y$};
    % 未来光锥
    \draw[thick] (0, 0) -- (-1.55, 1.85);
    \draw[thick] (0, 0) -- (1.55, 1.85);
    \draw[thick] (0, 1.85) ellipse (1.55 and 0.46);
    \draw[thick, dashed] (-1.55, 1.85) arc (180:360:1.55 and 0.46);
    % 过去光锥
    \draw[thick, dashed] (0, 0) -- (-1.32, -1.55);
    \draw[thick, dashed] (0, 0) -- (1.32, -1.55);
    \draw[thick, dashed] (0, -1.55) ellipse (1.32 and 0.40);
    % 三个向量
    \draw[->, very thick] (0, 0) -- (0, 2.15) node[right, pos=0.85] {$V$};
    \draw[->, very thick] (0, 0) -- (2.45, 0.62) node[above right] {$W$};
    \draw[->, very thick] (0, 0) -- (1.30, 1.55) node[above right, pos=0.95] {$Z$};
    % 顶点
    \fill (0, 0) circle (1.6pt) node[below left] {$p$};
\end{tikzpicture}%
}

% 2. 图 9.2: (1+1) 维 Minkowski 空间中的容许区域 Omega_{theta,theta'}
%    初始超曲面 t=0, 末超曲面 t=theta' 与 |x| = 1 - theta^{-1} t
\newcommand{\drawFigNineTwo}{%
\begin{tikzpicture}[scale=1.0, >=Stealth, line join=round]
    % 坐标轴
    \draw[->] (0, -0.60) -- (0, 2.95) node[above] {$t$};
    \draw[->] (-2.90, 0) -- (3.45, 0) node[right] {$x$};
    % 容许区域
    \fill[gray!18] (-2.2, 0) -- (2.2, 0) -- (0.75, 1.45) -- (-0.75, 1.45) -- cycle;
    % 初始超曲面 t = 0
    \draw[very thick] (-2.2, 0) -- (2.2, 0);
    % 末超曲面 t = theta'
    \draw[very thick] (-0.75, 1.45) -- (0.75, 1.45);
    % 两条类光边界 |x| = 1 - theta^{-1} t
    \draw[very thick] (2.2, 0) -- (0, 2.2);
    \draw[very thick] (-2.2, 0) -- (0, 2.2);
    \draw[dashed] (0, 2.2) -- (0, 2.85);
    % 引线
    \draw[densely dotted, gray] (1.05, 2.14) -- (0.33, 1.87);
    \draw[densely dotted, gray] (1.58, 1.50) -- (0.80, 1.46);
    % 标注
    \node[below] at (-1.60, -0.06) {\small $t = 0$};
    \node[right] at (1.62, 1.55) {\small $t = \theta'$};
    \node at (2.05, 2.28) {\small $|x| = 1 - \theta^{-1}t$};
    \node at (-1.05, 0.50) {\small $\Omega_{\theta, \theta'}$};
\end{tikzpicture}%
}

% 3. 图 9.3: 容许时空区域 Omega 及其修改 Omega_{delta_I=eta_I} (delta_I, eta > 0 时扩大)
\newcommand{\drawFigNineThree}{%
\begin{tikzpicture}[scale=1.0, >=Stealth, line join=round]
    \draw[->] (0, -1.05) -- (0, 2.55) node[above] {$t$};
    \draw[->] (-3.05, 0) -- (3.15, 0) node[right] {$x$};
    % 扩大后的区域 (虚线)
    \draw[dashed, thick] (-2.40, -0.30) -- (2.40, -0.30) -- (0.33, 2.22) -- (-0.33, 2.22) -- cycle;
    % 原区域 Omega (实线)
    \fill[gray!15] (-2.05, 0) -- (2.05, 0) -- (0.13, 1.95) -- (-0.13, 1.95) -- cycle;
    \draw[very thick] (-2.05, 0) -- (2.05, 0) -- (0.13, 1.95) -- (-0.13, 1.95) -- cycle;
    % 引线
    \draw[densely dotted, gray] (-1.62, -0.60) -- (-1.50, -0.02);
    \draw[densely dotted, gray] (1.34, 2.00) -- (0.46, 1.97);
    \draw[densely dotted, gray] (0.92, -0.62) -- (0.66, -0.14);
    % 标注
    \node at (-0.88, 0.62) {$\Omega$};
    \node[above] at (-1.70, -0.98) {\small 初始超曲面};
    \node[right] at (1.38, 2.02) {\small 末超曲面};
    \node[right] at (0.95, -0.80) {\small $\Omega_{\delta_I = \eta_I}$};
\end{tikzpicture}%
}

% 4. 图 9.4: 从 Omega_{delta,-eta} 出发, 终止于 Omega_{delta,-eta} \setminus Omega_{0,-eta} 的类光测地线
%    由于 delta_I = delta > 0 而 delta_F = -eta < 0, 壳层位于初始超曲面 t = 0 的下侧
\newcommand{\drawFigNineFour}{%
\begin{tikzpicture}[scale=1.0, >=Stealth, line join=round]
    \draw[->] (0, -1.40) -- (0, 1.95) node[above] {$t$};
    \draw[->] (-2.95, 0) -- (3.15, 0) node[right] {$x$};
    % 外层区域 Omega_{delta,-eta}
    \fill[gray!10] (-2.60, -1.15) -- (2.60, -1.15) -- (1.17, 1.60) -- (-1.17, 1.60) -- cycle;
    % 壳层 Omega_{delta,-eta} \setminus Omega_{0,-eta}
    \fill[gray!35] (-2.60, -1.15) -- (2.60, -1.15) -- (2.00, 0) -- (-2.00, 0) -- cycle;
    % 边界
    \draw[very thick] (-2.60, -1.15) -- (2.60, -1.15) -- (1.17, 1.60)
        -- (-1.17, 1.60) -- cycle;
    \draw[very thick] (-2.00, 0) -- (2.00, 0);
    % 过去指向的类光测地线 (斜率 -1)
    \draw[very thick, ->] (-0.40, 1.20) -- (1.75, -0.95);
    \fill (-0.40, 1.20) circle (1.8pt);
    \fill (1.75, -0.95) circle (1.8pt);
    % 引线
    \draw[densely dotted, gray] (-1.42, 1.15) -- (-0.46, 1.19);
    \draw[densely dotted, gray] (1.87, 1.32) -- (1.24, 1.06);
    % 标注
    \node[left] at (-1.45, 1.15) {\small $\pi(\gamma(0))$};
    \node[below left] at (1.28, -0.70) {\footnotesize $\pi(\gamma(T))$};
    \node[right] at (1.90, 1.32) {\small $\Omega_{\delta, -\eta}$};
    \node at (-0.95, 0.45) {\small $\Omega_{0, -\eta}$};
    \node at (-1.15, -0.55) {\footnotesize $\Omega_{\delta, -\eta} \setminus \Omega_{0, -\eta}$};
\end{tikzpicture}%
}

% 5. 图 9.5: 估计 (9.116) 的图示
%    Omega_{-delta,0} = Omega \cap {t >= delta} 为缩小后的区域, 壳层 Omega \setminus Omega_{-delta,0} 位于初始超曲面附近
\newcommand{\drawFigNineFive}{%
\begin{tikzpicture}[scale=1.0, >=Stealth, line join=round]
    \draw[->] (0, -0.55) -- (0, 2.35) node[above] {$t$};
    \draw[->] (-2.85, 0) -- (3.20, 0) node[right] {$x$};
    % 整个 Omega
    \fill[gray!12] (-2.35, 0) -- (2.35, 0) -- (0.15, 2.10) -- (-0.15, 2.10) -- cycle;
    % 壳层 Omega \setminus Omega_{-delta,0} (初始超曲面附近的带状区域)
    \fill[gray!40] (-2.35, 0) -- (2.35, 0) -- (1.82, 0.55) -- (-1.82, 0.55) -- cycle;
    \draw[very thick] (-2.35, 0) -- (2.35, 0) -- (0.15, 2.10) -- (-0.15, 2.10) -- cycle;
    % Omega_{-delta,0} 的初始超曲面 t = delta
    \draw[thick, dashed] (-1.82, 0.55) -- (1.82, 0.55);
    % 引线
    \draw[densely dotted, gray] (0.90, 1.74) -- (0.45, 1.62);
    % 标注
    \node at (-0.78, 0.88) {\footnotesize $\Omega_{-\delta, 0}$};
    \node[right] at (0.95, 1.78) {\small $\Omega$};
    \node at (1.05, 0.22) {\footnotesize $\Omega \setminus \Omega_{-\delta, 0}$};
\end{tikzpicture}%
}

% 6. 图 9.6: 纤维 T*_{(0,z)}M' 及其线性子空间与仿射子空间
\newcommand{\drawFigNineSix}{%
\begin{tikzpicture}[scale=1.0, >=Stealth, line join=round]
    % 纤维 (以平行四边形示意)
    \fill[gray!8] (-3.1, -1.35) -- (2.9, -1.35) -- (3.6, 1.35) -- (-2.4, 1.35) -- cycle;
    \draw[thick] (-3.1, -1.35) -- (2.9, -1.35) -- (3.6, 1.35) -- (-2.4, 1.35) -- cycle;
    \node[above left] at (-2.4, 1.35) {\small $T_{(0,z)}^* M'$};
    % 线性子空间 T*_z M
    \draw[very thick] (-2.85, -0.55) -- (3.2, -0.55);
    \node[below] at (2.9, -0.60) {\small $T_z^* M$};
    % 仿射子空间 ds + T*_z M
    \draw[very thick, dashed] (-2.75, 0.60) -- (3.3, 0.60);
    \node[above] at (2.85, 0.62) {\small $\mathrm{d}s + T_z^* M$};
    % 特征集 Char(P')
    \draw[very thick] plot[smooth, tension=0.75] coordinates {(-2.2, 1.05) (-0.9, 0.72) (0.35, 0.30) (1.55, -0.28) (2.5, -0.85)};
    \node[right] at (2.52, -0.92) {\small $\Char(P')$};
    % 交点
    \fill (-0.05, 0.60) circle (2pt);
    \fill (0.85, -0.55) circle (2pt);
    \node[above right] at (-0.03, 0.62) {\small $\Char_h(P)$};
\end{tikzpicture}%
}
